% Ch. 32 : le DAG d'entraînement de Listify
\documentclass[border=6pt]{standalone}
\usepackage{coursfig}
\begin{document}
\begin{tikzpicture}[x=1mm,y=1mm]
  \tikzset{t/.style={box=#1, text width=30mm, minimum height=14mm}}
  \node[t=Plum] (ext) at (0,0) {\textbf{extraire}\sub{export du jour}};
  \node[t=Plum] (val) at (42,0) {\textbf{valider}\sub{schéma, catégories}};
  \node[t=Enc] (ent) at (84,0) {\textbf{entraîner}\sub{run MLflow}};
  \node[t=Ocre] (dec) at (126,0) {\textbf{décider}\sub{branchement}};
  \node[t=Sea] (pro) at (172,12) {\textbf{promouvoir}\sub{alias champion}};
  \node[t=Brick] (ale) at (172,-14) {\textbf{alerter}\sub{le candidat est refusé}};
  \foreach \a/\b in {ext/val, val/ent, ent/dec} \draw[flow] (\a) -- (\b);
  \draw[flow=Sea] (dec.east) -- ++(8,0) |- (pro.west);
  \draw[flow=Brick] (dec.east) -- ++(8,0) |- (ale.west);
  \node[elabel, anchor=south, text=Sea] at (140,14) {précision $\ge$ seuil};
  \node[elabel, anchor=north, text=Brick] at (149,-16) {sinon};
  \node[note, anchor=north west, text width=80mm] at (0,-22) {chaque tâche est un processus distinct : elle peut échouer, être relancée, s'exécuter sur une autre machine};
\end{tikzpicture}
\end{document}
